% ====================================================================
\chapter{Functions and Polynomials}\label{ch:func}
% ====================================================================

\begin{chapterintro}
Economic models are built from relationships: quantity demanded depends on price, cost depends on output, utility
depends on consumption. The mathematical object that captures a relationship of this kind is a \emph{function}. This
chapter gives the informal and formal definitions of a function, then studies the simplest important family of numerical
functions, the polynomials, and proves that a polynomial can be factorised through each of its roots.
\end{chapterintro}

\begin{prerequisites}
Sets and set-builder notation (\cref{ch:sets}); summation notation (\cref{ch:sum}); the strong form of induction
(\cref{sec:strong-ind}); factorising quadratics through their roots (\cref{sec:quadratic}).
\end{prerequisites}

\section{The concept of a function}\label{sec:function}

\begin{definition}[Function (informal)][def:function-informal]
A \term{function} $f\colon D\to R$ is a rule that associates each element of a set $D$ (the \term{domain}) with an
element of a set $R$ (the \term{codomain}). \src{L2, slide 16}
\end{definition}

\begin{note}[Notation]
In the slides the codomain is called $R$, an ordinary capital letter for an arbitrary set. Do not confuse it with $\R$,
the set of real numbers. When the codomain is the real numbers, we write $f\colon D\to\R$.
\end{note}

Two features of the definition matter. \emph{Every} element of the domain is associated with something: a function is
defined everywhere on its domain. And each element is associated with exactly \emph{one} element of the codomain: a
function never gives two outputs for one input. Several inputs may share an output, and some elements of the codomain may
not be outputs at all.

\begin{definition}[Function (formal)][def:function-formal]
Given sets $D$ and $R$, a \term{function} $f$ is a set of ordered pairs $(x,y)$, with $x\in D$ and $y\in R$, such that for
every $x\in D$ there is a unique $y\in R$ with $(x,y)\in f$. In this case we write $y=f(x)$. \src{L2, slide 16}
\end{definition}

The formal definition identifies a function with its \emph{graph}: the set of all input--output pairs. It shows that a
function is a set (\cref{ch:sets}) with a special property, and it removes the vague word ``rule''. ``For every $x\in D$
there is a unique $y$'' combines a $\forall$, an $\exists$ and a uniqueness condition (\cref{sec:quantifiers}).

\begin{example}[Two functions]
\begin{enumerate}
  \item Let $D$ be a set of physical objects, $R$ a set of masses, and $g(x)$ the mass of object $x$. For instance,
  $g(\text{Jupiter})\approx1.898\times10^{27}$\,kg.
  \item Let $D=R=\R$ and $h(x)=x^2$. For instance $h(10)=10^2=100$ and $h(-1)=(-1)^2=1$.
\end{enumerate}
\src{L2, slide 16}
In (b), the inputs $1$ and $-1$ share the output $1$, which is allowed; and negative numbers in the codomain are never
outputs, which is also allowed.
\end{example}

\begin{definition}[Equal functions][def:function-equal]
Two functions $f$ and $g$ are \term{equal} if they have the same domain $D$ and the same codomain $R$, and for any
$x\in D$, $f(x)=g(x)$. \src{L2, slide 16}
\end{definition}

\begin{example}[Same formula, different functions]
(a) $f\colon\R\to\R$, $f(x)=x^2$, and $g\colon[0,\infty)\to\R$, $g(x)=x^2$, are \emph{not} equal: their domains differ.
(b) $p\colon\R\setminus\{1\}\to\R$, $p(x)=\frac{x^2-1}{x-1}$, and $q\colon\R\setminus\{1\}\to\R$, $q(x)=x+1$, \emph{are}
equal: on their common domain $\frac{x^2-1}{x-1}=\frac{(x-1)(x+1)}{x-1}=x+1$. But $p$ is not equal to
$r\colon\R\to\R$, $r(x)=x+1$, because $p$ is not defined at $x=1$.
\end{example}

\begin{definition}[Numerical function][def:numerical-function]
When $D,R\subseteq\R$, we call $f\colon D\to R$ a \term{numerical function}. \src{L2, slide 16}
\end{definition}

Almost every function in ECN115 is numerical, or (in Topic~6) a function of several real variables.

\begin{external}[Functions in economics]
\begin{itemize}
  \item A \emph{demand function} $Q=D(p)$ gives the quantity demanded at each price $p\ge0$, for instance $D(p)=100-2p$
  on the domain $[0,50]$.
  \item A \emph{cost function} $C(q)$ gives the total cost of producing $q\ge0$ units.
  \item A \emph{production function} $Y=F(K,L)$ gives output from capital $K$ and labour $L$; it is a function of two
  variables.
\end{itemize}
The domain is part of the model: a negative price or quantity usually has no meaning, so the domain is restricted to
non-negative values.
\end{external}

\section{Polynomials}\label{sec:polynomials}

\begin{definition}[Polynomial][def:polynomial]
Given $n\in\{0\}\cup\N$ and numbers $a_0,a_1,\dots,a_n\in\R$, define a function $f\colon\R\to\R$ by
\[
f(x)=a_0+a_1x+a_2x^2+\dots+a_nx^n=\sum_{k=0}^{n}a_kx^k.
\]
\begin{itemize}
  \item The function $f$ is called a \term{polynomial}.
  \item The numbers $a_0,a_1,\dots,a_n$ are called the \term{coefficients} of $f$.
  \item If at least one coefficient is not zero, the largest number $k$ such that $a_k\ne0$ is called the
  \term{degree} of $f$.
  \item The functions $a_0,\ a_1x,\ a_2x^2,\ \dots,\ a_nx^n$ are called the \term{monomials} (terms) of $f$.
  \item If $f(\bar x)=0$, we call $\bar x$ a \term{root} of the polynomial $f$.
\end{itemize}
\src{L2, slide 17}
\end{definition}

\begin{example}[Degrees and roots]
\begin{enumerate}
  \item $f(x)=3x^2-5x+2$ has degree $2$ and roots $\tfrac23$ and $1$ (\cref{ex:quad-ineq}).
  \item $f(x)=7$ has degree $0$ (the only coefficient is $a_0=7\ne0$) and no roots.
  \item $f(x)=0$ (all coefficients zero) has no degree under the definition, and every real number is a root.
  \item $f(x)=0x^3+2x+1$ has degree $1$, not $3$: the degree is the largest power with a \emph{non-zero} coefficient.
\end{enumerate}
\end{example}

Polynomials of degree $1$, $2$ and $3$ are called \emph{linear}, \emph{quadratic} and \emph{cubic}. A non-zero
constant polynomial has degree $0$ and no root, which matters in the proof below.

\section{Factorising a polynomial through a root}\label{sec:factor-theorem}

For a quadratic with roots $x_1,x_2$ we found $ax^2+bx+c=a(x-x_1)(x-x_2)$. The next theorem shows that every root of every
polynomial produces a linear factor in the same way.

\begin{theorem}[Polynomial factorisation via roots (Theorem L2.1)][thm:factor]
Suppose a polynomial $p$ of degree $n\ge1$, $p(x)=a_nx^n+\dots+a_1x+a_0$, has a root $\bar x\in\R$; that is,
$p(\bar x)=0$. Then there is a unique polynomial $q$ of degree $n-1$ such that
\[
p(x)=(x-\bar x)\cdot q(x)\qquad\text{for all }x\in\R. \tag*{\srcin{L2, slide 18}}
\]
\end{theorem}

\begin{core}[Root $\Leftrightarrow$ linear factor]
$\bar x$ is a root of $p$ exactly when $x-\bar x$ is a factor of $p$. (If $p(x)=(x-\bar x)q(x)$ then clearly $p(\bar x)=0$;
the theorem gives the converse.) Finding roots and finding linear factors are the same problem.
\end{core}

\emph{Proof (by induction on the degree $n$).} \src{L2, slides 18--19}

\emph{Induction base, $n=1$.} Then $p(x)=a_1x+a_0$ with $a_1\ne0$. The polynomial $p$ has the unique root
$\bar x=-a_0/a_1$, and
\[
p(x)=a_1\Big(x+\frac{a_0}{a_1}\Big)=(x-\bar x)\cdot a_1=(x-\bar x)\cdot q(x),
\]
where $q(x)=a_1$ is a polynomial of degree $1-1=0$. This $q$ is unique, since $(x-\bar x)\cdot q(x)=a_1x+a_0$ implies
$q(x)=a_1$ (compare the coefficients of $x$).

\emph{Induction step.} Consider an arbitrary $n>1$ and assume that the statement of the theorem is true for all degrees
$k$ with $1\le k<n$ (this is the alternative formulation of induction, \cref{sec:strong-ind}). Let $p$ have degree $n$
and root $\bar x$.

\emph{Case 1: $p(x)=(x-\bar x)\cdot a_nx^{n-1}$.} Then $q(x)=a_nx^{n-1}$, a polynomial of degree $n-1$, does the job.

\emph{Case 2: $p(x)\ne a_nx^{n-1}(x-\bar x)$.} Consider the difference
\[
s(x)=p(x)-a_nx^{n-1}(x-\bar x)=(a_{n-1}+a_n\bar x)\,x^{n-1}+a_{n-2}x^{n-2}+\dots+a_0.
\]
The $x^n$ terms cancel, so $s$ is a polynomial whose degree is at most $n-1$, and $s$ is not the zero polynomial (this is
what Case 2 means). Moreover
\[
s(\bar x)=p(\bar x)-a_n\bar x^{\,n-1}(\bar x-\bar x)=0-0=0,
\]
so $\bar x$ is a root of $s$. A non-zero polynomial of degree $0$ is a non-zero constant and has no root, so the degree
$k$ of $s$ satisfies $1\le k\le n-1$. By the induction hypothesis, $s(x)=(x-\bar x)\cdot r(x)$ for a (unique) polynomial
$r$ of degree $k-1\le n-2$. In this case
\[
p(x)=s(x)+a_nx^{n-1}(x-\bar x)=(x-\bar x)\cdot\big(a_nx^{n-1}+r(x)\big).
\]
The polynomial $q(x)=a_nx^{n-1}+r(x)$ has degree exactly $n-1$, because $a_n\ne0$ and $r$ has degree at most $n-2$. This
completes the existence part of the step. \hfill$\square$

\begin{sourcenote}[Uniqueness in the induction step]
The slides label the final $q$ ``a unique $q(x)$ of degree $n-1$'' but do not show why no other polynomial works. The
uniqueness of $r$ from the induction hypothesis does not by itself rule out a different factorisation of $p$. The missing
argument is supplied below; it is short and uses only the existence part.
\end{sourcenote}

\emph{Proof of uniqueness.} Suppose $p(x)=(x-\bar x)q_1(x)=(x-\bar x)q_2(x)$ for all $x$, where $q_1,q_2$ are polynomials.
Let $d=q_1-q_2$, a polynomial with $(x-\bar x)d(x)=0$ for every $x$. For every $x\ne\bar x$ we may divide by $x-\bar x\ne0$,
so $d(x)=0$: every real number other than $\bar x$ is a root of $d$. If $d$ were a non-zero polynomial, of degree $m$ say,
then applying the existence part repeatedly would factor out one linear factor for each of $m+1$ distinct roots, leaving a
non-zero polynomial of negative degree, which is impossible; equivalently, a non-zero polynomial of degree $m$ has at most
$m$ roots (\cref{prop:at-most-n}). So $d$ is the zero polynomial and $q_1=q_2$. \hfill$\square$

\begin{proposition}[At most $n$ roots][prop:at-most-n]
A polynomial of degree $n\ge0$ has at most $n$ distinct roots.
\end{proposition}
\emph{Proof.} By induction on $n$. A polynomial of degree $0$ is a non-zero constant and has no roots. Suppose the claim holds
for degree $n-1$, and let $p$ have degree $n\ge1$. If $p$ has no root we are done. Otherwise, let $\bar x$ be a root; by
\cref{thm:factor} (existence part), $p(x)=(x-\bar x)q(x)$ with $q$ of degree $n-1$. If $y\ne\bar x$ is another root, then
$0=p(y)=(y-\bar x)q(y)$ with $y-\bar x\ne0$, so $q(y)=0$ (\cref{prop:zero-product}). Hence every root of $p$ other than
$\bar x$ is a root of $q$, of which there are at most $n-1$; so $p$ has at most $n$ roots. \hfill$\square$

\subsection{Finding the factor in practice}

\begin{example}[Factorising a cubic][ex:cubic]
Factorise $p(x)=x^3-6x^2+11x-6$ completely.
\solution
Try small integers: $p(1)=1-6+11-6=0$, so $1$ is a root and, by \cref{thm:factor}, $p(x)=(x-1)q(x)$ with $q$ quadratic.
Write $q(x)=ax^2+bx+c$ and compare coefficients in $(x-1)(ax^2+bx+c)=ax^3+(b-a)x^2+(c-b)x-c$:
\[
a=1,\qquad b-a=-6,\qquad c-b=11,\qquad -c=-6,
\]
giving $a=1$, $b=-5$, $c=6$ (and the third equation, $6-(-5)=11$, checks). So $p(x)=(x-1)(x^2-5x+6)=(x-1)(x-2)(x-3)$.
\end{example}

\begin{external}[Where to look for integer roots]
If a polynomial has integer coefficients and leading coefficient $1$, any integer root must divide the constant term. In
\cref{ex:cubic} the candidates were $\pm1,\pm2,\pm3,\pm6$.
\end{external}

\begin{example}[Using the factorisation to solve an inequality]
Solve $x^3-6x^2+11x-6>0$.
\solution
By \cref{ex:cubic} the expression is $(x-1)(x-2)(x-3)$. Critical points $1,2,3$. For $x>3$ all three factors are positive;
crossing each critical point changes the sign of exactly one factor. So the product is positive on $(3,\infty)$, negative
on $(2,3)$, positive on $(1,2)$ and negative on $(-\infty,1)$. Solution: $x\in(1,2)\cup(3,\infty)$.
\end{example}

\begin{example}[Break-even points of a cubic profit function]
A firm's profit is $\pi(q)=-q^3+9q^2-23q+15$ (in thousands of pounds, with $q$ in thousands of units). Find all outputs at
which the firm breaks even ($\pi=0$), and where profit is positive.
\solution
$\pi(1)=-1+9-23+15=0$, so $q=1$ is a root. Dividing as in \cref{ex:cubic}: $\pi(q)=-(q-1)(q^2-8q+15)=-(q-1)(q-3)(q-5)$. The
firm breaks even at $q=1,3,5$. The product $(q-1)(q-3)(q-5)$ is negative on $(3,5)$ and on $(-\infty,1)$, so $\pi(q)>0$ for
$q\in[0,1)\cup(3,5)$ (restricting to $q\ge0$).
\end{example}

% --------------------------------------------------------------------
\section{Chapter review}
% --------------------------------------------------------------------
\subsection*{Summary}
A function $f\colon D\to R$ associates each element of its domain $D$ with exactly one element of its codomain $R$;
formally it is a set of ordered pairs $(x,y)$ with a unique $y$ for each $x$. Two functions are equal only if their domains,
codomains and values all agree. Numerical functions have $D,R\subseteq\R$. A polynomial is
$f(x)=\sum_{k=0}^na_kx^k$; its degree is the largest $k$ with $a_k\ne0$, and a root is a value where it vanishes.
Theorem L2.1 says that a root $\bar x$ of a polynomial $p$ of degree $n\ge1$ yields a unique factorisation
$p(x)=(x-\bar x)q(x)$ with $q$ of degree $n-1$; consequently a polynomial of degree $n$ has at most $n$ roots.

\subsection*{Key definitions}
\begin{itemize}
  \item \textbf{Function:} a rule (formally, a set of ordered pairs) giving each $x\in D$ a unique $f(x)\in R$.
  \item \textbf{Domain, codomain:} the sets $D$ and $R$ in $f\colon D\to R$.
  \item \textbf{Equal functions:} same domain, same codomain, same value at every point.
  \item \textbf{Numerical function:} $D,R\subseteq\R$.
  \item \textbf{Polynomial, coefficients, degree, monomials, root:} as in \cref{def:polynomial}.
\end{itemize}

\subsection*{Key results}
\begin{gather*}
p(\bar x)=0\ \Longleftrightarrow\ p(x)=(x-\bar x)\,q(x)\text{ with }\deg q=\deg p-1;\\
\text{a polynomial of degree }n\text{ has at most }n\text{ roots.}
\end{gather*}

\subsection*{Connections}
Functions are sets of ordered pairs (\cref{ch:sets}), and polynomials are written with $\sum$ (\cref{ch:sum}). The proof of
Theorem L2.1 uses the strong form of induction (\cref{sec:strong-ind}), and the theorem generalises the factorisation of
quadratics used to solve inequalities (\cref{ch:ineq}). Sequences, studied next, are functions with domain $\N$
(\cref{ch:limits}). Topic~2 continues the study of functions.

\subsection*{Common mistakes}
\begin{itemize}
  \item Regarding two formulas as the same function when their domains differ.
  \item Reading the degree from the largest power written rather than the largest power with a non-zero coefficient.
  \item Confusing the codomain $R$ in $f\colon D\to R$ with the real numbers $\R$.
  \item Forgetting that the factor theorem requires $\bar x$ to be a root.
\end{itemize}

\subsection*{Exam checklist}
\begin{checklist}
  \item Give the informal and formal definitions of a function, and the definition of equal functions.
  \item Define polynomial, coefficient, degree, monomial and root.
  \item State Theorem L2.1 and outline its proof by induction.
  \item Factorise a polynomial given (or after finding) a root, by comparing coefficients.
\end{checklist}

% --------------------------------------------------------------------
\section{Practice questions}
% --------------------------------------------------------------------
\pq[basic] Which of the following define a function from $\R$ to $\R$? (a) $f(x)=x^3$; (b) $f(x)=\sqrt x$;
(c) $f(x)=\frac1{x^2+1}$; (d) $f(x)$ is ``a number whose square is $x^2$''.

\pq[basic] State the degree and the roots of: (a) $2x-8$; (b) $x^2+1$; (c) $(x-1)^2(x+4)$; (d) $0\cdot x^2+5$.

\pq[conceptual] Are the functions $f\colon\R\to\R$, $f(x)=|x|$, and $g\colon\R\to\R$, $g(x)=\sqrt{x^2}$, equal? Are
$f\colon\R\to\R$, $f(x)=x$, and $h\colon\R\to[0,\infty)$, $h(x)=|x|$, equal? Justify.

\pq[mathematical] Show that $2$ is a root of $p(x)=x^3-3x^2-4x+12$, and factorise $p$ completely.

\pq[mathematical] Use \cref{prop:at-most-n} to show that if two polynomials of degree at most $n$ agree at $n+1$ distinct
points, they are equal.

\pq[application] A demand function is $D(p)=120-3p$. State a sensible domain, find the price at which demand is zero, and
explain why $D$ cannot be defined with domain $\R$ in an economic model.

\pq[multi-step] Find all real $x$ satisfying $x^3-2x^2-5x+6\le0$.

\pq[exam-style] (a) State Theorem L2.1 on the factorisation of a polynomial via its roots. (b) Explain where the proof uses
the induction hypothesis and why the alternative (strong) formulation of induction is convenient. (c) Use the theorem to
factorise $p(x)=2x^3-3x^2-11x+6$, given that $p(3)=0$.

% --------------------------------------------------------------------
\section{Solutions to practice questions}
% --------------------------------------------------------------------
\pqsol{9.1}
(a) Yes. (b) No: $\sqrt x$ is not defined for $x<0$, so not every element of the domain $\R$ has an output. (c) Yes: the
denominator is never $0$. (d) No: for $x\ne0$ there are two such numbers ($x$ and $-x$), so the output is not unique.

\pqsol{9.2}
(a) Degree $1$, root $4$. (b) Degree $2$, no real roots. (c) Degree $3$, roots $1$ and $-4$. (d) Degree $0$ (the polynomial
is the constant $5$), no roots.

\pqsol{9.3}
$f$ and $g$ are equal: same domain and codomain, and $\sqrt{x^2}=|x|$ for every real $x$. $f$ and $h$ are not equal: their
codomains differ, and in any case $f(-1)=-1\ne1=h(-1)$.

\pqsol{9.4}
$p(2)=8-12-8+12=0$. Then $p(x)=(x-2)(x^2+bx+c)$; comparing, $b-2=-3$, so $b=-1$, and $-2c=12$, so $c=-6$. Thus
$p(x)=(x-2)(x^2-x-6)=(x-2)(x-3)(x+2)$.

\pqsol{9.5}
If $f$ and $g$ have degree at most $n$ and agree at $n+1$ points, then $d=f-g$ is a polynomial of degree at most $n$ (or
the zero polynomial) with at least $n+1$ roots. By \cref{prop:at-most-n} a non-zero polynomial of degree at most $n$ has at
most $n$ roots, so $d$ is the zero polynomial and $f=g$.

\pqsol{9.6}
Quantities and prices cannot be negative: $p\ge0$ and $D(p)\ge0\Leftrightarrow p\le40$. A sensible domain is $[0,40]$, with
demand zero at $p=40$. With domain $\R$, the formula would assign negative quantities to prices above $40$ and to negative
prices, which have no economic meaning.

\pqsol{9.7}
$p(1)=1-2-5+6=0$, so $p(x)=(x-1)(x^2-x-6)=(x-1)(x-3)(x+2)$. Critical points $-2,1,3$. For $x>3$ the product is positive;
it is negative on $(1,3)$, positive on $(-2,1)$, negative on $(-\infty,-2)$. Including the roots:
$x\in(-\infty,-2]\cup[1,3]$.

\pqsol{9.8}
(a) If a polynomial $p$ of degree $n\ge1$ has a root $\bar x$, there is a unique polynomial $q$ of degree $n-1$ with
$p(x)=(x-\bar x)q(x)$.
(b) In Case 2 the auxiliary polynomial $s(x)=p(x)-a_nx^{n-1}(x-\bar x)$ has the root $\bar x$ and degree $k$ with
$1\le k\le n-1$, but $k$ need not equal $n-1$; the hypothesis is applied to $s$, of degree $k$. The strong form, which
assumes the theorem for \emph{all} degrees below $n$, covers every possible $k$; the ordinary form would cover only
$k=n-1$.
(c) $p(x)=(x-3)(2x^2+bx+c)$; comparing, $b-6=-3$ so $b=3$, and $-3c=6$ so $c=-2$. Hence $p(x)=(x-3)(2x^2+3x-2)$, and
$2x^2+3x-2=(2x-1)(x+2)$, so $p(x)=(x-3)(2x-1)(x+2)$.
